% ============================================================
% English Abstract
% ============================================================

With the rapid advancement of the electronics industry, component selection and supply chain risk management have become critical bottlenecks affecting product development efficiency and reliability. Traditional manual selection workflows rely heavily on engineer experience, are time-consuming and error-prone, and lack systematic supply chain risk assessment. This paper presents an intelligent component selection and risk assessment Agent system for the eZ-PLM electronic component management platform, built upon Large Language Models (LLM) and Retrieval-Augmented Generation (RAG) technology.

The system adopts a five-stage pipeline architecture: natural language requirement parsing, multi-dimensional knowledge retrieval, intelligent scoring and ranking, risk identification, and report generation. A LangChain ReAct Agent paradigm is further introduced to enable autonomous multi-tool orchestration and multi-turn dialogue. Key innovations include: (1) a dual-mode requirement understanding mechanism combining LLM-prioritized parsing with rule-based fallback, supporting Chinese and English natural language input across Buck, Boost, and LDO topologies; (2) a ChromaDB-based engineering knowledge vector database that injects design formulas, standards, and best practices into the selection decision chain via semantic retrieval; (3) a five-dimensional weighted scoring model (parameter matching, supply chain stability, application fit, design maturity, and evidence confidence) with dual-mode switching between rule-only and LLM-enhanced scoring; (4) three standardized Intermediate Representation schemas—PartIR, RiskIR, and TopologyIR—enabling data interoperability with the eZ-PLM platform.

Experimental results demonstrate 100\% pass rate across 10 DC-DC buck converter selection test cases. RAG knowledge enhancement significantly improves the engineering rigor of topology analysis reports, which now automatically include inductor calculation formulas, thermal performance estimation, and PCB layout guidance. The ReAct Agent mode supports multi-turn interactions including domestic alternative inquiries, with tool calls controlled within 2–4 per turn. The generated BOM lists, risk assessment reports, and topology analysis documents are ready for engineering review, validating the practical value of AI Agents in electronic engineering.
